\documentclass[12pt,a4paper]{article}
\usepackage{fontspec}
\usepackage{xeCJK}
\setmainfont{Times New Roman}
\setCJKmainfont{Songti SC}[BoldFont=Heiti SC, AutoFakeBold=2.2]
\setCJKsansfont{Heiti SC}[AutoFakeBold=3.2]
\usepackage{geometry}
\geometry{a4paper,left=2.15cm,right=2.15cm,top=2.05cm,bottom=2.25cm}
\usepackage{setspace}
\setstretch{1.5}
\usepackage{enumitem}
\setlist{nosep,leftmargin=1.7em,itemsep=0.18em,topsep=0.28em}
\usepackage{fancyhdr}
\pagestyle{fancy}
\fancyhf{}
\renewcommand{\headrulewidth}{0pt}
\fancyfoot[C]{24}
\setlength{\parindent}{2em}
\begin{document}
\section*{附录}
\subsection*{附录 A\quad 支撑材料文件列表}
\begin{itemize}
\item 问题一：\texttt{问题一/问题一代码final\_修正版.py}（半平面交 + 旋转卡壳 + Welzl + 覆盖判定 + 仿真检验）、\texttt{问题一\_表格与图像.py}、\texttt{problem1\_localize.py}、\texttt{figs/}（示向度、交会定位、角形区域及验证图）。
\item 问题二：\texttt{问题二/code/problem2\_model.py}、\texttt{problem2\_paper.py}、\texttt{problem1\_localize.py}、\texttt{test\_problem2.py}、\texttt{final\_refine.py}、\texttt{results/problem2\_summary.json}、\texttt{figs/}（三张候选区域图）。
\item 问题三：\texttt{问题三/code/}（\texttt{problem3\_belief\_struct.py}、\texttt{problem3\_belief.py}、\texttt{problem3\_route\_solver.py}、\texttt{run\_official\_e15.py}、\texttt{problem3\_official\_robot.py} 等上场源程序）、\texttt{正式日志/}（\texttt{formal-p3-1-S589-ME5P-AVEG-7FVF.jlog}、\texttt{formal-p3-2-CRNF-M248-R8CD-3QE2.jlog}、\texttt{formal-p3-3-G3WW-HEAT-UQXS-8K7V.jlog} 及表 1）、\texttt{本地结果/}（fullopt 1 万局 \texttt{summary.json}/\texttt{trials.csv}，E15 本地 1000 局摘要）。
\item 问题四：\texttt{问题四/code/}（\texttt{problem4\_belief\_n19.py}、\texttt{run\_official\_n19.py} 及依赖源程序）、\texttt{正式测试成绩.txt}（三次官方成绩；未附 \texttt{formal-p4-*.jlog}）、\texttt{本地结果/}（N19 本地 1000 局 \texttt{summary.json}/\texttt{trials.csv}）。
\item AI 工具使用详情：\texttt{AI工具使用详情.pdf}。
\end{itemize}

\subsection*{附录 B\quad 源程序代码}
完整、可运行的源程序代码见支撑材料，本附录不全文粘贴。
\end{document}
